\documentclass[11pt]{article}

\usepackage[margin=1.2in]{geometry}
\usepackage{amsmath,amssymb}
\usepackage{microtype}
\usepackage{xcolor}
\usepackage{mdframed}
\usepackage[colorlinks=true,linkcolor=blue!50!black,urlcolor=blue!50!black]{hyperref}

\newmdenv[backgroundcolor=blue!4,linecolor=blue!40,linewidth=1pt,
  skipabove=8pt,skipbelow=8pt]{keyidea}

\title{The Collatz Conjecture at the Critical Line:\\
       \Large A Friendly Guide}
\author{Companion notes to the paper of the same title\thanks{Written in
collaboration with Claude (Anthropic). These notes explain the ideas; the
paper and the Lean files contain the proofs. Nothing here is a proof, and
nothing here claims the conjecture --- it remains open.}}
\date{June 12, 2026}

\begin{document}
\maketitle

\section{The game}

Take any positive whole number. If it is even, halve it. If it is odd,
multiply by three and add one. Repeat.
\[
7 \to 22 \to 11 \to 34 \to 17 \to 52 \to 26 \to 13 \to 40 \to 20 \to 10
\to 5 \to 16 \to 8 \to 4 \to 2 \to 1.
\]
The \textbf{Collatz conjecture} says: no matter which number you start
with, you always end at 1. Computers have checked every starting number
up to about $2.4$ sextillion ($2^{71}$). Nobody has proved it always
happens, and nobody has found a number that escapes. This guide explains
a collection of theorems --- each checked line-by-line by a computer
proof assistant called Lean, so there is no ``trust me'' anywhere ---
that map out exactly \emph{why} the problem is hard and \emph{how close}
mathematics can currently get.

One bookkeeping convenience first. When $n$ is odd, $3n+1$ is always
even, so the next move is always a halving. It saves ink to fuse them:
the \emph{Terras map} sends odd $n$ to $(3n+1)/2$ and even $n$ to $n/2$.
Every step now includes exactly one halving, which makes size accounting
clean. ``Odd step'' (roughly multiply by $1.5$) versus ``even step''
(multiply by $0.5$) is the whole tug-of-war.

\section{Why it should be true --- and why that's not a proof}

Here is the back-of-envelope argument every mathematician meets first.
If the odd/even choices behaved like fair coin flips, a typical step
would multiply your number by about
$\sqrt{1.5 \times 0.5} = \sqrt{0.75} \approx 0.87$ --- shrinking. So a
``random'' orbit drifts down and eventually hits 1. The catch: orbits
are not random. They are completely determined by the starting number.
The conjecture is a claim about \emph{every individual} number, and
averages say nothing about the worst case. The gap between ``almost
every number behaves'' and ``every number behaves'' is, as we will see,
precisely where the whole problem lives.

\begin{keyidea}
\textbf{The critical ratio.} Suppose the first $T$ steps contain $j$ odd
steps. Up to small corrections your number has been multiplied by
$3^j / 2^T$. That exceeds 1 exactly when $j/T$ exceeds
\[
\frac{\log 2}{\log 3} = 0.6309\ldots
\]
So: an orbit climbs while more than about $63\%$ of its steps are odd,
and falls otherwise. Everything in this story orbits around that one
number, the \emph{critical density}.
\end{keyidea}

\section{The tempting proof strategy, and the theorem that kills it}

Because any finite range can be checked by computer, the most natural
plan is:
\begin{enumerate}
\item Sort numbers into finitely many bins --- say, by remainder when
dividing by $2^k$.
\item Check, bin by bin, that every number is guaranteed to drop below
its starting value within some fixed number of steps.
\item Conclude every large number eventually falls into the
already-verified range, and you are done.
\end{enumerate}
Step 3 is fine. Step 2 \emph{feels} like a finite computation, and
countless attempted proofs rest on some version of it. The first main theorem says step 2 is not just unproven
but \textbf{impossible}: no binning, no clever weighting of bins, no
fixed step budget can ever work.

\textbf{The villain.} Consider $n = 2^L - 1$: in binary, a string of
$L$ ones. It is odd, so apply $(3n+1)/2$. A short calculation gives
$3 \cdot 2^{L-1} - 1$ --- whose binary form ends in $L-1$ ones. Still
odd. Apply again: ends in $L-2$ ones. The number climbs by a factor of
about $1.5$ per step for $L$ \emph{consecutive} steps: $100\%$ odd
density, far above critical, for as long as you arranged ones in
advance.

Why does this defeat \emph{every} binning argument? A bin ``knows''
only a bounded amount about a number --- its remainder, the state of
some finite gadget, a bounded score. But $2^L - 1$ exists for every
$L$. Whatever the bound, choose $L$ so large that the all-ones number
climbs straight through any fixed step budget; the bookkeeping that
was supposed to certify descent gets exhausted. (In fancier language:
$-1$ behaves like a fixed point of the odd step --- check:
$(3 \cdot (-1) + 1)/2 = -1$ --- and $2^L - 1$ imitates $-1$ for $L$
steps. The fancier language is what the formal proof uses, but the
binary picture is the same fact.)

\begin{keyidea}
\textbf{The no-go theorem.} Any descent certificate whose memory of $n$
is bounded --- remainder classes for any modulus, any finite automaton,
any bounded weighting, any fixed window of steps --- is defeated by the
numbers $2^L - 1$. This closes off an entire genre of proof attempts at
once, including weighted and amortized versions.
\end{keyidea}

A concrete sample, verified in Lean down to the witness: the plausible
claim ``every $n > 300$ drops below itself within 53 steps'' is false,
and the smallest counterexample is $447$ --- binary
$110111111$, whose six trailing ones are a pocket-sized version of the
villain.

\section{The information budget: why the first $\log_2 n$ steps are
special}

Here is the cleanest way to understand both the obstruction and the
positive results. When you halve a number, you consume its last binary
digit. A number with $\log_2 n$ digits can therefore only ``answer
questions'' about its first $\log_2 n$ halvings; beyond that, the digits
are spent and the orbit's choices are fresh information you never had.

This intuition is a theorem (Terras, 1976, formalized here):
\begin{itemize}
\item the pattern of odd/even choices in the first $k$ steps is
\emph{exactly determined} by the remainder of $n$ upon division by
$2^k$ --- and conversely, the pattern determines the remainder;
\item consequently \textbf{every} odd/even pattern of length $k$
is realized by exactly one remainder class. All behaviors occur. The
all-ones villain is just one cell of a complete table;
\item counting cells: the number of remainder classes whose first $k$
steps contain exactly $j$ odd ones is the binomial coefficient
$\binom{k}{j}$ --- Pascal's triangle, the same count as coin-flip
sequences.
\end{itemize}

\begin{keyidea}
\textbf{Why the problem is genuinely two-faced.} Every finite behavior
occurs (so worst-case reasoning is hopeless --- that is the no-go
theorem), yet behaviors are distributed like coin flips (so typical
numbers fall --- that is where the ``almost all'' results come from).
\end{keyidea}

\section{The corridor: exact bookkeeping}

The heuristic ``multiplied by $3^j/2^T$'' can be made exact, with the
error terms under complete control. Two inequalities, proved by
induction with nothing but integer arithmetic, sandwich the truth:
\[
3^j(n+1) - 2^j \;\le\; 2^T \cdot (\text{value after } T \text{ steps})
\;\le\; 3^j (n + 2^{T-j}) - 2^T .
\]
Both become exact equalities on the villain orbit, so neither can be
improved. From the sandwich follow the two walls of the corridor:

\textbf{Wall 1 (forced density).} If $n$ has not yet dropped below its
starting value after $T$ steps, then either its odd-step share already
exceeds the critical $63.09\%$, or it has used up more than $\log_2 n$
halvings --- its entire information budget. Inside the budget, survival
\emph{requires} supercritical density. No exceptions.

\textbf{Wall 2 (the wall is touched).} Conversely, for every window
length $B$ there really are arbitrarily large numbers that survive all
$B$ steps with odd density only a hair above $63.09\%$ --- one can
prescribe the pattern ``$j$ odd steps then $s$ even ones'' and, by the
pattern-realization theorem, manufacture actual integers that follow it.
The achievable survival densities form a strip of width about $1/B$
sitting exactly on the critical line, inhabited at every scale.

And a reduction theorem ties the knot: the full conjecture is
\emph{equivalent} to the statement that every $n \ge 2$ eventually drops
below itself. So:

\begin{keyidea}
\textbf{What the conjecture actually is.} A number escapes to infinity
only by keeping its odd-step share above $63.09\%$ forever ---
threading a strip that narrows like $1/B$ at every scale
simultaneously. The strip is never empty at any finite scale, and no
bounded-memory argument can decide who threads it. That tension
\emph{is} the Collatz problem.
\end{keyidea}

\section{Almost every number falls (and it's elementary)}

Combine the walls with Pascal's triangle. To survive $k$ steps a number
needs a supercritical pattern, but among the $2^k$ remainder classes
only $\sum_{j \gtrsim 0.63k} \binom{k}{j}$ have one --- an exponentially
small minority, since coin flips concentrate near $50\%$. Conclusion
(Terras' theorem, here in fully finite form): among the numbers between
$2^k$ and $2^{k+1}$, the fraction that fails to drop within $k$ steps is
at most about $2^{-0.045k}$ --- vanishing exponentially.

Two touches make the formal version pleasant. The binomial tail is
tamed by a trick worth showing off: multiply each tail term by
$2^{j}/2^{m} \geq 1$ and use
\[
2^m \sum_{j \ge m} \binom{k}{j} \;\le\; \sum_{j} \binom{k}{j} 2^j
= (1+2)^k = 3^k,
\]
the binomial theorem --- no calculus, no entropy, nothing beyond
algebra. And the threshold ``$63.09\%$'' is replaced by the fraction
$17/27 = 0.6296\ldots$, justified by the single inequality
$3^{17} \le 2^{27}$. The price of staying this elementary is a slightly
generous constant; the prize is that a computer can certify every step.

\section{Why there are no small loops}

A number could also fail by entering a loop (other than the familiar
$1 \to 2 \to 1$). Going once around a loop of $T$ steps with $j$ odd
ones multiplies by almost exactly $3^j/2^T$ --- and must return you to
where you started, so $3^j/2^T \approx 1$, i.e.\ $3^j \approx 2^T$. But
a power of 3 never \emph{equals} a power of 2, and near-misses are
governed by how well the irrational number $\log_2 3$ can be
approximated by fractions $T/j$ --- which is ``badly,'' most of the
time. The sandwich inequalities turn this into a clean bound: every
member of a loop is smaller than $3^j 2^{T-j} / (2^T - 3^j)$.

Feed in the computer-verified fact that everything below $2^{71}$
reaches 1, and a finite check (performed by Lean's proof kernel itself,
all $17{,}000$ or so cases) yields: \textbf{no loop with 183 or fewer
halvings exists}. The first case the method cannot kill is $T = 184$,
$j = 116$, where $3^{116}$ approximates $2^{184}$ just well enough ---
the wall is a genuine number-theoretic fact, not a software limitation.

\section{Beyond the theorem: the experimental frontier}

The no-go theorem forbids \emph{bounded}-memory certificates. What
about scores that read \emph{all} the binary digits of $n$ --- for
instance $V(n) = n \cdot c^{(\text{number of ones in } n)}$? These
genuinely escape the theorem's hypotheses, and on the villain they
work: the all-ones orbit loses ones as it climbs, so for $c > 9/4$ the
score falls even while the number rises. The papers report a systematic
machine search (with the search's own counterexample-finder) over such
digit-reading scores. Every candidate died, and the autopsies found
three recurring killers:
\begin{enumerate}
\item \textbf{Injection}: start from a number whose digits make the
score artificially tiny (lots of zeros); the dynamics restores typical
digit statistics and the score balloons --- even while the number
itself shrinks.
\item \textbf{Global rewrite}: an odd step multiplies the whole binary
string by 3; digit counts get scrambled non-locally, so no digit-based
score is steered by the digits it read.
\item \textbf{Shift release}: a single halving shifts the whole string,
instantly cashing in whatever ``stored credit'' a position-sensitive
score had accumulated --- the best candidate found (it effectively
predicted the villain's climbs in advance) died this way in one step.
\end{enumerate}
The pattern of the three killers points somewhere specific: binary
(base-2) statistics are destroyed by multiplication by 3, and base-3
statistics would be destroyed by halving. Anything that survives both
must live on base 2 and base 3 \emph{simultaneously} --- and the
interaction of doubling and tripling is a famously deep area of
mathematics (Furstenberg's ``$\times 2, \times 3$'' problems). The
experiments, in other words, independently rediscover where the experts
always suspected the difficulty lives.

\section{What ``machine-verified'' means here}

Every theorem above is written in Lean 4, a proof assistant whose small
trusted kernel checks each logical step; ``the computer checked it''
means checked the \emph{proof}, not merely tested examples. The files
compile with no unproven placeholders, and an audit confirms no axioms
beyond Lean's standard foundations. The one statement taken on faith is
clearly labeled as a hypothesis: that computers have verified the
conjecture below $2^{71}$ (a published computation, not formalized).

\section{Summary in five sentences}

Numbers fall on average but the conjecture is about every number.
Any proof that remembers only finitely much about $n$ is provably
impossible, because all-ones numbers climb past any budget. Each number
controls only its first $\log_2 n$ steps, behaviors within that budget
are distributed like coin flips, and survival demands an exponentially
rare $63\%$-odd pattern --- so almost every number falls, provably and
elementarily. Escape would mean riding the $63.09\%$ critical line at
every scale forever, through a strip that narrows like $1/B$ yet is
never empty; loops would need $3^j \approx 2^T$ and are ruled out up to
183 halvings. What remains open is exactly the gap between ``almost
every'' and ``every'' --- and that gap is now fenced, measured, and
machine-checked on both sides.

\end{document}
